\section{Workload Characterization}\label{sec:workload}

Before any policy comparison can be meaningful, we need to know what the
underlying workloads look like. This section collects both raw-trace statistics
and a picture of actual traffic for Congress.gov and CourtListener, then
explains why replay-Zipf overlay is required for either trace to yield a
well-posed eviction-policy benchmark.

\subsection*{Trace collection}

\textbf{Congress.gov API.} We collected $20{,}692$ requests against the
Congress.gov v3 REST API at a conservative $1.2\,$s base delay with
exponential backoff and randomized jitter, sweeping the
\texttt{/bill/}, \texttt{/amendment/}, and \texttt{/vote/} endpoint families.
Requests were client-generated random queries against the key space (bill,
amendment, and vote identifiers drawn from the cataloged archive), not captured
user traffic; this is a deliberate design choice since no public capture of
production Congress.gov queries exists.

\textbf{CourtListener REST v4.} We chose CourtListener REST v4 because
PACER's \$500--\$2000 paywall made a $20$\,K-request trace infeasible. The
resulting trace contains $20{,}000$ requests at $0.8\,$s base pacing with a
$0.4\,$s random jitter, mixed $80/20$ between metadata endpoints
(\texttt{/dockets/}, \texttt{/clusters/}, \texttt{/courts/}) and full-text
opinion fetches. Collection ran overnight ($\approx 9$ hours wall-clock) under
a single authenticated token. CourtListener serves judicial documents that
get cited across briefs, dockets refreshed during active cases, and cluster
identifiers that appear across multiple search results --- so the trace
contains genuine structural locality that survives into the request log.

\subsection*{Trace characteristics}

Table~\ref{tab:workload} reports the ten summary statistics that drive every
subsequent mechanism argument.

\begin{table}[h]
\centering
\footnotesize
\renewcommand{\arraystretch}{0.92}
\begin{tabularx}{\textwidth}{@{}lXX@{}}
\toprule
\textbf{Metric} & \textbf{Congress} & \textbf{Court} \\
\midrule
Total requests        & 20{,}692       & 20{,}000 \\
Unique objects        & 18{,}970       & 15{,}018 \\
Zipf $\alpha$ (MLE)   & 0.231          & 1.028 \\
OHW ratio (10\% win)  & 0.989          & 1.000 \\
Mean size (bytes)     & 751.9          & 3{,}144.6 \\
Median size (bytes)   & 231            & 1{,}381 \\
p95 size (bytes)      & 2{,}698        & 6{,}221 \\
Max size (bytes)      & 6{,}700        & 462{,}490 \\
Working set (bytes)   & 14{,}490{,}463 & 59{,}650{,}083 \\
\bottomrule
\end{tabularx}
\caption{Workload characterization --- Congress.gov vs.\ CourtListener raw
traces. \textbf{Note the $335\times$ max/median ratio on Court: a single
$462$\,KB document dominates the byte-miss-ratio analysis in
\S\ref{sec:bytemrc}.}}
\label{tab:workload}
\end{table}

\subsection*{Workload structure visualization}

\begin{figure}[h]
\centering
\begin{minipage}{0.48\textwidth}
\includegraphics[width=\textwidth]{../results/congress/figures/workload.pdf}
\end{minipage}\hfill
\begin{minipage}{0.48\textwidth}
\includegraphics[width=\textwidth]{../results/court/figures/workload.pdf}
\end{minipage}
\caption{Workload visualization --- Congress (left) vs.\ Court (right). The
per-request timeline and size distribution grounds the
$\alpha_{\text{raw}}=0.231$ vs.\ $1.028$ contrast numerically summarized in
Table~\ref{tab:workload}.}
\label{fig:workload_viz}
\end{figure}

Figure~\ref{fig:workload_viz} shows Congress's near-uniform access pattern
side-by-side with Court's long-tail shape. The Congress trace looks essentially
random over the key space; the Court trace concentrates on a small hot set of
judicial documents that get cited and re-fetched during active briefing.
These raw shapes --- not the replay-Zipf overlay we build on top of them ---
drive the mechanism divergence we explain in \S\ref{sec:mechanism}.

\subsection*{Why the raw $\alpha$ differs}

The $4.4\times$ gap between $\alpha_{\text{raw}}=0.231$ and $1.028$ is not a
measurement artifact; it reflects what each API is used for. Congress.gov API
calls are client-generated random queries over the bill, amendment, and vote
key space --- no user naturally re-requests bill \#1234 immediately after
bill \#5678, and there is no front-end aggregator that would force a
re-request pattern. The observed $\alpha=0.23$ matches ``essentially uniform
random over the key space,'' as confirmed by the unique-key ratio
($18{,}970/20{,}692 = 91.7\%$) and the near-total one-hit-wonder fraction
($0.989$ of keys are requested only once in the window covered by the trace).

CourtListener's REST v4, by contrast, serves judicial documents in a
real access pattern: opinions get cited in briefs, dockets get refreshed
during active cases, the same cluster-ID appears across multiple search
results from related queries. Real structural locality survives in the
trace, producing $\alpha_{\text{raw}}=1.028$ and a unique-key ratio
($15{,}018/20{,}000 = 75.1\%$) materially lower than Congress's.

\subsection*{Why we still run replay-Zipf on both}

Both traces are overwhelmingly one-hit-wonder dominated (OHW ratios
$0.989$ and $1.000$), so running the policies on the raw captures yields
near-indistinguishable results --- every admission is a miss, every eviction
victim is a key that would never be re-accessed anyway. Neither trace, in
its captured form, gives meaningful policy-comparison results. We therefore
overlay a replay-Zipf popularity distribution: the real object keys and sizes
are preserved, but the popularity ranks are resampled under a controlled
$\alpha$ sweep. This is what makes cross-workload comparison well-posed ---
differences in policy behavior trace back to the raw trace's object
population and size distribution, not to an artifact of the request-interval
pattern (which neither public-records API exposes in production form
anyway).
